\documentclass[11pt,a4paper]{article}
\usepackage[T1]{fontenc}
\usepackage[utf8]{inputenc}
\usepackage{amsmath,amsthm,mathtools}
\usepackage{newtxtext,newtxmath}
\usepackage[a4paper,margin=25mm,headheight=15pt]{geometry}
\usepackage{microtype,graphicx,booktabs,array,enumitem}
\usepackage[dvipsnames]{xcolor}
\usepackage[most]{tcolorbox}
\usepackage{fancyhdr,titlesec,xurl}
\usepackage[pdfencoding=auto,psdextra]{hyperref}
\hypersetup{colorlinks=true,linkcolor=MidnightBlue,citecolor=MidnightBlue,
 urlcolor=MidnightBlue,pdftitle={Partitions into Catalan Numbers: All-Orders Asymptotics, Periodic Corrections, and Inverse Growth},
 pdfauthor={AI-assisted research report prepared for Vladimir Reshetnikov},bookmarksnumbered=true}
\usepackage[nameinlink,noabbrev]{cleveref}
\numberwithin{equation}{section}
\newtheorem{theorem}{Theorem}[section]
\newtheorem{proposition}[theorem]{Proposition}
\newtheorem{lemma}[theorem]{Lemma}
\newtheorem{corollary}[theorem]{Corollary}
\theoremstyle{definition}
\newtheorem{definition}[theorem]{Definition}
\theoremstyle{remark}
\newtheorem{remark}[theorem]{Remark}
\newcommand{\R}{\mathbb R}
\newcommand{\Z}{\mathbb Z}
\newcommand{\N}{\mathbb N}
\newcommand{\ii}{\mathrm i}
\newcommand{\e}{\mathrm e}
\newcommand{\dd}{\mathop{}\!\mathrm d}
\newcommand{\HH}{\mathcal H}
\newcommand{\PP}{\mathbb P}
\newcommand{\coeff}[2]{\left[#1\right]#2}
\newcommand{\aset}{\mathcal A}
\DeclareMathOperator{\Rea}{Re}
\DeclareMathOperator{\Var}{Var}
\DeclareMathOperator{\Li}{Li}
\setlength{\parskip}{4pt}
\setlength{\parindent}{1.2em}
\setlength{\emergencystretch}{2em}
\setcounter{tocdepth}{2}
\titleformat{\section}{\large\bfseries\color{MidnightBlue}}{\thesection}{0.7em}{}
\titleformat{\subsection}{\normalsize\bfseries}{\thesubsection}{0.7em}{}
\pagestyle{fancy}
\fancyhf{}
\fancyhead[L]{\small Partitions into Catalan numbers}
\fancyhead[R]{\small OEIS A033552}
\fancyfoot[C]{\thepage}
\renewcommand{\headrulewidth}{0.3pt}
% Collection edition (batch 73O1): notes added when the report was written
% into ProveIt's research-report collection.
\newcommand{\writenote}[1]{\par\smallskip\noindent
 \textit{[Write note, batch 73O1, 1 October 2026.]}\ #1\par\smallskip}
\newtcolorbox{resultbox}{enhanced,breakable,colback=blue!2!white,
 colframe=MidnightBlue!65!black,boxrule=.5pt,arc=1mm,left=9pt,right=9pt,top=7pt,bottom=7pt}
\title{\vspace{-1.1cm}\textbf{Partitions into Catalan Numbers}\\[5pt]
\Large All-Orders Asymptotics, Periodic Corrections,\\ and Inverse Growth}
\author{A proof-focused research report for the ProveIt project}
\date{October 1, 2026}
\begin{document}
\maketitle
\thispagestyle{empty}
\begin{abstract}
Let $p(n)$ count ordinary integer partitions whose allowed part sizes are
$1,2,5,14,42,\ldots$, the distinct Catalan numbers. This is OEIS A033552.
We develop a complete expansion in inverse powers of a logarithmic size
parameter, including a nonconstant periodic amplitude, an explicit first
correction, and an inverse expansion for the threshold at which $p(n)$ reaches
a prescribed value. The argument combines an all-orders coefficient saddle
estimate, a Barnes $G$ product identity, and an exponentially convergent
Fourier series involving $\Gamma(s)\zeta(1+s)$. The contour estimate controls
all noncentral arcs; the expansion is therefore more than a formal saddle
calculation. A further theorem gives a phase-dependent limiting distribution
for the largest allowed-part index in a uniformly chosen partition.
Exact computations through $n=10^6$ and independent high-precision checks
accompany the proofs. The inspected OEIS entry contains no asymptotic formula;
we do not claim that it records an explicit asymptotic conjecture, or that an
exhaustive literature-priority search has been completed. The sequence-specific
formulas and proofs below are the substantive output, not a claim that the
classical saddle or periodic-partition methods themselves are new.
\end{abstract}

\begin{resultbox}
\textbf{Principal formula.} Put $\alpha=\log 4$ and let $\mu$ be the large
solution of $n=4^\mu/\sqrt{\pi\mu}$. Then
\[
 p(n)=\e^{C+\Psi_\alpha(\mu)}\mu^{23/8}
 \exp\!\left(\frac\alpha2\mu^2-\frac{1+3\alpha}{2}\mu\right)
 \left(1+\frac{D_\alpha(\mu)}\mu+O(\mu^{-2})\right).
\]
Here $C$ is explicit and $\Psi_\alpha,D_\alpha$ are smooth, one-periodic
functions given below. Every subsequent inverse-power coefficient is
constructively determined. Replacing $\Psi_\alpha$ by a constant does
\emph{not} give an asymptotic equivalent.
\end{resultbox}

\writenote{This is the collection edition of the report, written into
ProveIt's research-report collection from manuscript~39 of batch~73O1
(provenance in \cref{ctp:app:edition}). Every label carries the prefix
\texttt{ctp:}; \texttt{requirements.txt} now lives in \texttt{data/} and
\texttt{build.sh} in \texttt{code/}. Text added in the write is marked like
this note; everything else is the delivered text. As the abstract says,
the OEIS entry states \emph{no} asymptotic conjecture: this report
answers no posed question, it supplies the asymptotics the entry lacks.}

\clearpage
{\small\setlength{\parskip}{0pt}\tableofcontents}
\clearpage

\section{The problem, its status, and the contribution}
\label{ctp:sec:intro}

\subsection{The sequence and the normalization}
For $k\geq1$ define
\begin{equation}\label{ctp:eq:catalan}
 c_k=\frac1{k+1}\binom{2k}{k},\qquad
 F(z)=\prod_{k\geq1}(1-z^{c_k})^{-1},\qquad
 p(n)=[z^n]F(z).
\end{equation}
Thus the allowed part sizes are $c_1=1,c_2=2,c_3=5,c_4=14,\ldots$.
The initial terms are
\[
 1,1,2,2,3,4,5,6,7,8,10,11,13,14,17,19,22,24,27,30,\ldots.
\]
The index starts at $n=0$. It is important not to include both $c_0$ and
$c_1$: doing so would create two colors of the unit part and a different
sequence. The definition and generating function agree with OEIS
A033552~\cite{oeis}.

The inspected entry supplies a generating function, a recursive computation,
and a reference to Pak's enumerative-complexity survey, but no asymptotic
formula. Pak explicitly discusses partitions into Catalan numbers among
restricted-partition counting problems that may be less tractable than some
classical examples~\cite[\S2.5]{pak}. That discussion concerns exact counting
complexity. An asymptotic expansion, however accurate to fixed order, does
not by itself settle that complexity question.

All logarithms in this article are natural. A periodic function is evaluated
on a real argument by periodic extension. All $O$-estimates concern the limit
of that argument, or of $n$, to positive infinity; their constants may depend
on a fixed truncation order but not on a fractional phase.

\subsection{What is established here, and what is not}
The principal contribution is a self-contained chain from the exact product
\eqref{ctp:eq:catalan} to the following results: an all-orders saddle expansion
with a contour proof; a completely normalized, phase-resolved explicit
expansion; a first correction containing derivatives of the periodic
amplitude; an inverse threshold expansion with error transport; and a
phase-dependent law for the largest part index. A general proposition
explains which coefficients depend only on the exponential and polynomial
growth of the allowed part sizes.

Periodicity in asymptotics of geometrically restricted partitions is a
classical phenomenon. Erd\H{o}s and Richmond discuss the de Bruijn--Mahler
phenomenon and its extensions~\cite{erdosrichmond}. Mellin analysis of
harmonic sums is developed systematically by Flajolet, Gourdon, and
Dumas~\cite{fgd}. We use that context, but prove the concrete estimates and
normalizations needed here rather than importing a vaguely matched theorem.
The explicit Catalan-part formulas below were derived for this report. A
limited search found no source giving this complete collection of formulas;
that observation is not a proof of publication priority.

\writenote{The phenomenon Erd\H{o}s and Richmond name goes back to de
Bruijn's treatment of Mahler's partition problem \cite{debruijn}, which
the delivered bibliography did not list: partitions into powers of a
fixed integer, whose logarithm carries a periodic function of a
logarithmic variable built from $\Gamma(s)\zeta(1+s)$ on the Mellin
lattice. The Catalan parts $c_k\sim4^k/(\sqrt\pi k^{3/2})$ are a
polynomially perturbed geometric sequence of that kind, and
\cref{ctp:lem:Fourier} is the corresponding Mellin computation; see also
the note after its proof.}

The manuscript contains conventional mathematical proofs and executable
checks. It is not a Lean formalization. The numerical checks corroborate,
but do not replace, the arguments. No computer search is used as a substitute
for a missing contour estimate or for a missing asymptotic remainder.

\subsection{Relation to the ProveIt repository}
The inspected repository snapshot is commit
\texttt{e6190a94552f119e9a2140e5fea8450a045a0837}. Its transseries documentation
separates asymptotic scales from inversion, distinguishes a continuous inverse
from an integer staircase inverse, and records qualifications on formalized
and unformalized results~\cite{proveit}. Those distinctions are useful here.
The companion discusses asymptotics of Catalan numbers themselves, whereas
our object is the partition sequence whose \emph{allowed parts} are Catalan
numbers. An initial repository code search for \texttt{A033552} returned no match;
that is a duplication check, not a mathematical premise.

The present article uses an exactly invertible dominant size coordinate and
then transports residual errors through inversion. Its periodic coefficients
also make a useful boundary case for a purely polynomial--logarithmic
transseries calculus: the coefficient algebra must retain a periodic phase.

\writenote{At the repository's current state the search result above
still holds: no report or formal development treats A033552 or
partitions into Catalan numbers. The pin quoted above is a commit; the
delivered \texttt{SOURCE\_NOTES.md} also names
\texttt{1f1981f682b2878bde51a6ad40c22777f362fc05}, the commit its first
repository request returned, and calls it ``a tree SHA, not a commit
SHA''. That is mistaken: it is a ProveIt commit of 1~October~2026 (the
root tree of \texttt{e6190a945} is the separate object
\texttt{8e77886c}). Neighbouring material is listed in
\cref{ctp:app:edition}. No statement of this report is formalized.}

\section{Main asymptotic and inverse formulas}
\label{ctp:sec:main}

\subsection{Constants and the periodic amplitude}
Set
\begin{equation}\label{ctp:eq:constants}
 \alpha=\log4,\qquad A=\frac{1+3\alpha}{2},\qquad
 \eta=\frac{23}{8},\qquad
 C=\frac14\log2+\frac34\log\pi+\log G(3/2),
\end{equation}
where $G$ is the Barnes function, normalized by $G(1)=1$ and
$G(z+1)=\Gamma(z)G(z)$~\cite{dlmfG}.
For the moment allow $\alpha>0$ to vary, and define
\begin{align}
 \kappa&=\frac{\pi^2}{12}-\frac{\gamma^2}{2}-\gamma_1,
       \label{ctp:eq:kappa}\\
 \Psi_\alpha(\theta)
 &=\frac\alpha{12}+\frac\kappa\alpha+
 \frac1\alpha\sum_{\ell\in\Z\setminus\{0\}}
 \Gamma\!\left(\frac{2\pi\ii\ell}{\alpha}\right)
 \zeta\!\left(1+\frac{2\pi\ii\ell}{\alpha}\right)
 \e^{2\pi\ii\ell\theta}.\label{ctp:eq:phase-fourier}
\end{align}
Here $\gamma$ is Euler's constant and $\gamma_1$ is the first Stieltjes
constant, with convention
$\zeta(1+s)=s^{-1}+\gamma-\gamma_1s+O(s^2)$.
The Fourier series and all of its derivatives converge absolutely.
Conjugate terms pair to make $\Psi_\alpha$ real.
A prime denotes differentiation with respect to $\theta$; a partial
$\alpha$-derivative holds $\theta$ fixed. The correction is
\begin{equation}\label{ctp:eq:D}
\boxed{\displaystyle
 D_\alpha(\theta)=\frac16-\frac1{2\alpha}
 -\frac32\partial_\alpha\Psi_\alpha(\theta)
 +\frac{17}{8\alpha}\Psi_\alpha'(\theta)
 -\frac{\Psi_\alpha'(\theta)^2+\Psi_\alpha''(\theta)}{2\alpha^2}.}
\end{equation}
Only after taking the partial derivative do we set $\alpha=\log4$.

Define the large real size coordinate
\begin{equation}\label{ctp:eq:mu}
 n=\frac{4^{\mu(n)}}{\sqrt{\pi\mu(n)}},\qquad
 \mu(n)=-\frac1{2\alpha}W_{-1}\!\left(-\frac{2\alpha}{\pi n^2}\right).
\end{equation}
The negative real branch $W_{-1}$ is intended. For sufficiently large $n$
this is the unique large solution and $\mu(n)\sim(\log n)/\alpha$.

\begin{theorem}[Explicit expansion, including the first correction]
\label{ctp:thm:explicit}
For the sequence \eqref{ctp:eq:catalan}, as $n\to\infty$ through integers,
\begin{align}
 \log p(n)={}&\frac\alpha2\mu^2-A\mu+\eta\log\mu+C+
 \Psi_\alpha(\mu)+\frac{D_\alpha(\mu)}\mu+O(\mu^{-2}).
 \label{ctp:eq:main-log}
\end{align}
Equivalently,
\begin{equation}\label{ctp:eq:main-multiplicative}
 p(n)=\e^{C+\Psi_\alpha(\mu)}\mu^\eta
 \e^{\alpha\mu^2/2-A\mu}
 \left(1+\frac{D_\alpha(\mu)}\mu+O(\mu^{-2})\right).
\end{equation}
For every fixed $J\geq0$ there are smooth, one-periodic functions
$L_1,\ldots,L_J$, with $L_1=D_\alpha$, such that
\begin{equation}\label{ctp:eq:all-orders-explicit}
 \log p(n)=\frac\alpha2\mu^2-A\mu+\eta\log\mu+C+
 \Psi_\alpha(\mu)+\sum_{j=1}^J\frac{L_j(\mu)}{\mu^j}
 +O_J(\mu^{-J-1}).
\end{equation}
Sections~\ref{ctp:sec:saddle}--\ref{ctp:sec:explicit-proof} prove this statement and
give a finite coefficient-construction procedure at each order.
\end{theorem}

The phrase ``all orders'' here means a remainder theorem for every fixed
truncation order. It does not assert convergence of the resulting infinite
inverse-power series, nor does it resolve exponentially smaller sectors.

\begin{corollary}[An unavoidable periodic factor]\label[corollary]{ctp:cor:oscillation}
There is no constant $B$ for which
\[
 p(n)\sim B\mu(n)^\eta\e^{\alpha\mu(n)^2/2-A\mu(n)}.
\]
More precisely, the limit points of the normalized sequence are exactly
\[
 \left[\exp\{C+\min_\theta\Psi_\alpha(\theta)\},\,
       \exp\{C+\max_\theta\Psi_\alpha(\theta)\}\right],
\]
an interval with distinct endpoints.
\end{corollary}
\begin{proof}
The gamma function has no zeros, and $\zeta(1+\ii t)\ne0$ for real
$t\ne0$~\cite{dlmfzeros}; hence the first Fourier coefficient in
\eqref{ctp:eq:phase-fourier} is nonzero. Thus $\Psi_\alpha$ is nonconstant.
For any fixed phase $\theta$, round
$4^{j+\theta}/\sqrt{\pi(j+\theta)}$ to the nearest integer $n_j$.
Implicit differentiation of \eqref{ctp:eq:mu} shows
$\mu(n_j)=j+\theta+o(1)$. Every phase therefore occurs as a subsequential
limit. Apply \eqref{ctp:eq:main-multiplicative} and continuity of $\Psi_\alpha$.
\end{proof}

The first Fourier harmonic has amplitude approximately
$0.00100279457386$. The modulation is small but does not tend to zero.
Its mean is approximately $0.64116608028709$. These are reproducible
floating-point evaluations, not interval-certified bounds.

\begin{corollary}[The first three logarithmic growth terms]
\label[corollary]{ctp:cor:logs}
With $X=\log n$,
\begin{equation}\label{ctp:eq:three-logs}
 \log p(n)=\frac{X^2}{2\alpha}+\frac{X\log X}{2\alpha}
 +\left(\frac{\log(\pi/\alpha)-1}{2\alpha}-\frac32\right)X
 +O\bigl((\log X)^2\bigr).
\end{equation}
\end{corollary}
\begin{proof}
The defining equation is
$X=\alpha\mu-\tfrac12\log\mu-\tfrac12\log\pi$. Hence
\[
 \alpha\mu=X+\tfrac12\log X+\tfrac12\log(\pi/\alpha)
 +O((\log X)/X).
\]
Substitute this into \eqref{ctp:eq:main-log} and collect terms of orders
$X^2$, $X\log X$, and $X$.
\end{proof}
The positive sign of $X\log X$ is a concrete effect of the factor
$k^{-3/2}$ in the growth of $c_k$. Replacing $c_k$ by $4^k$ too early loses
that information.

\subsection{The inverse threshold}
For real $y>1$ define the integer staircase
\begin{equation}\label{ctp:eq:threshold}
 N(y)=\min\{n\geq0:p(n)\geq y\}.
\end{equation}
Put
\begin{equation}\label{ctp:eq:R-s}
 Y=\log y,\qquad R=\sqrt{2\alpha Y},\qquad s=R/\alpha.
\end{equation}

\begin{theorem}[Inverse growth and its first correction]\label{ctp:thm:inverse}
As $y\to\infty$,
\begin{align}
 \log N(y)={}&R+A-\frac12\log s-\frac12\log\pi\notag\\
 &+\frac{\frac12(A^2-A)-\alpha\eta\log s
 -\alpha\{C+\Psi_\alpha(s+A/\alpha)\}}{R}
 +O\!\left(\frac{(\log R)^2}{R^2}\right).
 \label{ctp:eq:inverse}
\end{align}
In particular,
\begin{equation}\label{ctp:eq:inverse-leading}
 N(y)\sim
 8\sqrt{\frac{e\alpha}{\pi}}\,
 \frac{\exp\!\left(\sqrt{2\alpha\log y}\right)}
 {(2\alpha\log y)^{1/4}}.
\end{equation}
The inverse has an expansion to arbitrary fixed order in $R^{-1}$,
with coefficients polynomial in $\log R$ and smooth periodic functions of
$s+A/\alpha$. Section~\ref{ctp:sec:inverse-proof} proves the passage from a
smooth asymptotic inverse to the integer threshold.
\end{theorem}

\section{An all-orders coefficient saddle estimate}
\label{ctp:sec:saddle}

\subsection{Positive cumulants and the saddle}
For $t>0$ write
\begin{equation}\label{ctp:eq:H-K}
 H(t)=\log F(\e^{-t}),\qquad
 K_j(t)=(-1)^jH^{(j)}(t)
 =\sum_{k\geq1}c_k^j\sum_{\ell\geq1}\ell^{j-1}\e^{-\ell t c_k}
 \quad(j\geq1).
\end{equation}
Absolute convergence on each half-plane $\Rea t\geq t_0>0$ justifies
termwise differentiation. Since $K_1'(t)=-K_2(t)<0$, $K_1$ decreases
strictly from infinity to zero. Thus, for each $n>0$, there is a unique
$t_n>0$ satisfying
\begin{equation}\label{ctp:eq:saddle}
 K_1(t_n)=n.
\end{equation}
The real logarithmic interpolation of the part sizes is
\begin{equation}\label{ctp:eq:r}
 r(x)=\alpha x-\tfrac12\log\pi+
 \log\Gamma(x+\tfrac12)-\log\Gamma(x+2),\qquad x\geq1.
\end{equation}
It obeys $r(k)=\log c_k$. Since $r'(1)=1/2$ and
$r''(x)=\psi'(x+1/2)-\psi'(x+2)>0$, it is strictly increasing, with
$r'(x)\to\alpha$. Define
\begin{equation}\label{ctp:eq:m}
 u_n=\log(1/t_n),\qquad m_n=r^{-1}(u_n).
\end{equation}
The parameters $m_n$ and $\mu(n)$ are not identical; their small difference
is responsible for several cancellations in the explicit formula.

\begin{lemma}[Size of the scaled cumulants]\label[lemma]{ctp:lem:cumulants}
Let $t=\e^{-r(m)}$ and $m\to\infty$. For each fixed integer $j\geq1$,
\begin{equation}\label{ctp:eq:B}
 B_j(t):=t^jK_j(t)=(j-1)!\,m+O_j(1).
\end{equation}
In particular $nt_n\sim m_n$ and $m_n\sim(\log n)/\alpha$.
\end{lemma}
\begin{proof}
Set $f_j(v)=v^j\sum_{\ell\geq1}\ell^{j-1}\e^{-\ell v}$.
For $0<v\leq1$, $f_j(v)=(j-1)!+O_j(v)$, by the elementary Laurent
expansion of the geometric sum and its derivatives. For $v\geq1$,
$f_j(v)$ decays exponentially up to a fixed polynomial factor.
The exact ratio
\begin{equation}\label{ctp:eq:ratio}
 \frac{c_{k+1}}{c_k}=\frac{2(2k+1)}{k+2}\in[2,4)
\end{equation}
implies $\sum_{tc_k\leq1}tc_k=O(1)$ and a bounded sum of the
large-$v$ tails. There are exactly $\lfloor m\rfloor$ parts with
$tc_k\leq1$. Summing $f_j(tc_k)$ proves \eqref{ctp:eq:B}.
At the saddle, $nt=B_1=m+O(1)$, and
$\log n=r(m)+\log(m+O(1))=\alpha m+O(\log m)$.
\end{proof}

\begin{theorem}[All-orders saddle expansion]\label{ctp:thm:saddle}
For every fixed $J\geq0$, with all cumulants evaluated at $t=t_n$,
\begin{equation}\label{ctp:eq:saddle-expansion}
 p(n)=\frac{\exp(H(t)+nt)}{\sqrt{2\pi K_2(t)}}
 \left(1+\sum_{j=1}^J E_j(t)+O_J(m_n^{-J-1})\right).
\end{equation}
The correction $E_j$ is the finite sum
\begin{equation}\label{ctp:eq:Ej}
 E_j=\sum_{\substack{k_3,\ldots,k_{2j+2}\geq0\\
                   \sum_{r=3}^{2j+2}(r-2)k_r=2j}}
 (-1)^{R/2}(R-1)!!
 \prod_{r=3}^{2j+2}\frac1{k_r!}
 \left(\frac{K_r}{r!K_2^{r/2}}\right)^{k_r},
 \qquad R=\sum_{r=3}^{2j+2}rk_r.
\end{equation}
Here $R$ is even, and $E_j=O_j(m_n^{-j})$. In particular,
\begin{align}
 E_1={}&\frac{K_4}{8K_2^2}-\frac{5K_3^2}{24K_2^3},\label{ctp:eq:E1}\\
 E_2={}&-\frac{K_6}{48K_2^3}+\frac{7K_3K_5}{48K_2^4}
 +\frac{35K_4^2}{384K_2^4}-\frac{35K_3^2K_4}{64K_2^5}
 +\frac{385K_3^4}{1152K_2^6}.\label{ctp:eq:E2}
\end{align}
\end{theorem}

\subsection{The noncentral arcs: why the coefficient estimate is valid}
The Cauchy coefficient integral is
\begin{equation}\label{ctp:eq:cauchy}
 p(n)=\frac{\e^{H(t)+nt}}{2\pi}
 \int_{-\pi}^{\pi}
 \exp\{H(t-\ii\theta)-H(t)-\ii n\theta\}\dd\theta.
\end{equation}
Its normalized absolute value satisfies the exact identity
\begin{equation}\label{ctp:eq:modulus}
 \left|\frac{F(\e^{-t+\ii\theta})}{F(\e^{-t})}\right|
 =\prod_{k\geq1}
 \left(1+\frac{4\e^{-tc_k}\sin^2(c_k\theta/2)}
 {(1-\e^{-tc_k})^2}\right)^{-1/2}.
\end{equation}
All factors on the right are at most one.
Fix a requested error order. In the estimates below, $B>0$, the number of
fixed factors $K$, and a small $\delta>0$ are chosen in that order and then
held fixed as $m\to\infty$.

For $|\theta|\leq tm^B$, use every part with $c_k\leq1/(tm^B)$.
There are $m-O(\log m)$ such parts. For them,
$|c_k\theta|\leq1$, and elementary sine and exponential estimates give
\begin{equation}\label{ctp:eq:central-tail-bound}
 \left|\frac{F(\e^{-t+\ii\theta})}{F(\e^{-t})}\right|
 \leq\left(1+c(\theta/t)^2\right)^{-c'm}
\end{equation}
for fixed positive $c,c'$. Between
$t\log m/\sqrt m$ and $t$, integration of this bound gives
$O(t\exp\{-c''(\log m)^2\})$, after changing constants. Between $t$ and
$tm^B$, it gives $O(tm^B\e^{-c''m})$.
Both are smaller than $(t/\sqrt m)m^{-L}$ for every fixed $L$.

Next, for $tm^B\leq|\theta|\leq\delta$, select the first $K$ factors and
choose $\delta c_K\leq1$. Each selected factor in \eqref{ctp:eq:modulus} is
at most $C_Kt/|\theta|$. Hence
\begin{equation}\label{ctp:eq:middle-bound}
 \int_{tm^B\leq|\theta|\leq\delta}
 \left|\frac{F(\e^{-t+\ii\theta})}{F(\e^{-t})}\right|\dd\theta
 =O_K\left(t\,m^{-B(K-1)}\right).
\end{equation}
Taking $K$ sufficiently large makes this negligible to any specified
inverse-power order relative to the main width $t/\sqrt m$.

Finally consider $\delta\leq|\theta|\leq\pi$. The unit-part factor is
$O_\delta(t)$. The functions $\sin(2\theta/2)$ and
$\sin(5\theta/2)$ do not vanish simultaneously on this compact set:
a simultaneous zero would require $\theta\in2\pi\Z$ because
$\gcd(2,5)=1$. Therefore at least one of the part-$2$ and part-$5$
factors is uniformly $O_\delta(t)$. Thus
\begin{equation}\label{ctp:eq:outer-bound}
 \int_{\delta\leq|\theta|\leq\pi}
 \left|\frac{F(\e^{-t+\ii\theta})}{F(\e^{-t})}\right|\dd\theta
 =O_\delta(t^2).
\end{equation}
Since $t=\exp(-\alpha m+O(\log m))$, this is again negligible to every
inverse power of $m$ relative to $t/\sqrt m$.

These estimates eliminate possible competing root-of-unity arcs at the
precision claimed in \cref{ctp:thm:saddle}. They do not assert that all those
arcs vanish, or identify their exponentially smaller contributions.

\subsection{The central expansion and its remainder}
At the saddle the linear term in \eqref{ctp:eq:cauchy} cancels. With
$v=\theta\sqrt{K_2}$, the exponent becomes
\begin{equation}\label{ctp:eq:central-exponent}
 -\frac{v^2}{2}+
 \sum_{r\geq3}\frac{\ii^rK_r}{r!K_2^{r/2}}v^r.
\end{equation}
For each fixed $r$, \cref{ctp:lem:cumulants} gives
$K_r/K_2^{r/2}=O_r(m^{-(r-2)/2})$. The central interval
$|\theta|\leq t\log m/\sqrt m$ corresponds to $|v|\leq C\log m$.
Taylor estimates for $H(t-\ii\theta)$ follow directly from its absolutely
convergent sums: the absolute value of a derivative along this vertical
segment is bounded by the corresponding positive derivative sum at $t$.

Expand the exponential in \eqref{ctp:eq:central-exponent}, assigning weight
$r-2$ to the term of degree $r$. Retain total weight through $2J+1$.
The integrated terms of odd weight vanish because their total degree in
$v$ is odd. The weight-$2j$ terms have Gaussian integral
$(-1)^{R/2}(R-1)!!$ after division by $\sqrt{2\pi}$, giving
\eqref{ctp:eq:Ej}.

For completeness, the remainder estimate does not require retaining a
power of $\log m$. On $|v|\leq C\log m$, the absolute correction to the
Gaussian exponent is $O(m^{-1/2}(|v|^3+1))$, which is at most
$v^2/4+o(1)$ outside a fixed bounded interval. A Taylor remainder through
weight $2J+1$ is bounded by
\[
 C_Jm^{-J-1}(1+|v|^{M_J})\e^{-v^2/4}
\]
for a fixed integer $M_J$, after including the Gaussian factor and
increasing the Taylor order of $H$ to $2J+4$ if necessary. Its integral is
$O_J(m^{-J-1})$. The omitted Gaussian tails are superpolynomially small.
Combining this with \eqref{ctp:eq:central-tail-bound}--\eqref{ctp:eq:outer-bound}
proves \cref{ctp:thm:saddle}.

The leading values of the correction polynomials are useful checks:
\begin{equation}\label{ctp:eq:E-universal}
 E_1=-\frac1{12m}+O(m^{-2}),\qquad
 E_2=\frac1{288m^2}+O(m^{-3}),\qquad
 E_3=\frac{139}{51840m^3}+O(m^{-4}).
\end{equation}
They follow by substituting $B_r\sim(r-1)!m$ in \eqref{ctp:eq:Ej}.
At a finite $m$, using the exact cumulants is preferable to replacing each
$E_j$ by only its first universal term.

\subsection{A reusable sufficient hypothesis}
The same proof applies to an increasing sequence of positive integers
$a_1,a_2,\ldots$ when $a_1=1$, the ratios are eventually between two fixed
constants $1<q\leq Q<\infty$, and a finite subset of
$\{a_2,a_3,\ldots\}$ has greatest common divisor one. Replace $m$ by the
number of active parts $a_k\leq1/t$, up to an $O(1)$ difference. The finite
gcd condition is sufficient for the last compact arc estimate; it is not
claimed to be necessary. In particular, the proof as stated should not be
applied unchanged to pure powers of two, whose tail has a nontrivial gcd.

\section{The Catalan product and its moving periodic phase}
\label{ctp:sec:phase}

\subsection{The logarithm of a Catalan number and the Barnes product}
The standard gamma quotient in \eqref{ctp:eq:r} gives the complete expansion
\begin{equation}\label{ctp:eq:r-asymptotic}
 r(x)=\alpha x-\beta\log x+c+\sum_{j\geq1}b_jx^{-j},
 \quad
 \beta=\frac32,\quad c=-\frac12\log\pi,
 \quad b_1=-\frac98,\quad b_2=\frac12.
\end{equation}
The meaning is asymptotic, to every fixed order with differentiated
remainders. A convenient exact coefficient formula is
\begin{equation}\label{ctp:eq:bj}
 b_j=\frac{(-1)^{j+1}}{j(j+1)}
 \bigl(B_{j+1}(1/2)-B_{j+1}(2)\bigr),
\end{equation}
where $B_j(x)$ is a Bernoulli polynomial. For example,
$b_3=-21/64$. This follows by subtracting the shifted Stirling expansions
of the two log-gamma functions; it also gives an effective coefficient
source for every subsequent step.

For an integer $M\geq0$, put $S(M)=\sum_{k=1}^Mr(k)$. Repeated use of
$G(z+1)=\Gamma(z)G(z)$ gives the exact identity
\begin{equation}\label{ctp:eq:S-exact}
 S(M)=\frac\alpha2M(M+1)-\frac M2\log\pi+
 \log G(M+3/2)-\log G(M+3)-\log G(3/2).
\end{equation}
The Barnes expansion~\cite{dlmfG}, or Euler--Maclaurin summation followed
by the exact identity to fix the constant, yields
\begin{align}
 S(M)={}&\frac\alpha2M^2-\frac32M\log M+
 \left(\frac\alpha2-\frac12\log\pi+\frac32\right)M
 -\frac{15}{8}\log M+K \notag\\
 &-\frac{19}{16M}+\frac{65}{128M^2}+O(M^{-3}),
 \label{ctp:eq:S-asymptotic}\\
 K={}&-\frac34\log(2\pi)-\log G(3/2).\label{ctp:eq:K}
\end{align}
Subtracting the expansion at $M-1$ from the expansion at $M$ recovers
\eqref{ctp:eq:r-asymptotic}; this is a useful local check, but the Barnes
identity is what determines the constant $K$.

\subsection{Separating the bulk from a localized lattice sum}
Let
\begin{equation}\label{ctp:eq:g}
 g(v)=-\log(1-\e^{-\e^v})-(-v)_+,
 \qquad (-v)_+=\max\{-v,0\}.
\end{equation}
As $v\to-\infty$, $g(v)=O(\e^v)$; as $v\to+\infty$ it decays
superexponentially. It is continuous and piecewise analytic, with a jump
in its first derivative at zero. This mild kink will cancel against the
fractional-part polynomial.

For $m\geq1$, let $M=\lfloor m\rfloor$ and $\theta=m-M$. Define
\begin{equation}\label{ctp:eq:Hm}
 \HH(m)=H(\e^{-r(m)}).
\end{equation}
Then the following identity is exact:
\begin{equation}\label{ctp:eq:exact-decomposition}
 \HH(m)=Mr(m)-S(M)+\sum_{k\geq1}g(r(k)-r(m)).
\end{equation}
Indeed, the subtracted linear part $(-v)_+$ contributes $r(m)-r(k)$
precisely for $k\leq M$. If $m$ is an integer, the boundary term is zero,
so no convention ambiguity occurs.

Replacing $r(m+d)-r(m)$ by $\alpha d$ in the localized sum suggests
\begin{equation}\label{ctp:eq:phase-lattice}
 \Psi_\alpha(\theta)=\frac\alpha2\theta(1-\theta)
       +\sum_{j\in\Z}g(\alpha(j-\theta)),\qquad 0\leq\theta\leq1.
\end{equation}
The sum converges absolutely, exponentially fast on the left and faster
on the right. The polynomial is forced by expansion of $Mr(m)-S(M)$;
it is not an adjustable normalization.

\begin{lemma}[Fourier representation of the phase]\label[lemma]{ctp:lem:Fourier}
The lattice formula \eqref{ctp:eq:phase-lattice} equals
\eqref{ctp:eq:phase-fourier}. In particular it extends to a smooth,
nonconstant one-periodic function, and termwise derivatives of its Fourier
series of every fixed order are legitimate.
\end{lemma}
\begin{proof}
For $\Rea s>0$, change variables $y=\e^v$ and expand
$-\log(1-\e^{-y})=\sum_{\ell\geq1}\e^{-\ell y}/\ell$ to obtain
\[
 \int_\R g(v)\e^{sv}\dd v
 =\Gamma(s)\zeta(1+s)-\frac1{s^2}.
\]
The integral involving $g$ extends analytically to $\Rea s>-1$.
The apparent singularity on the right at zero is removable: the product
of the Laurent expansions of $\Gamma(s)$ and $\zeta(1+s)$ has leading
term $s^{-2}$, no $s^{-1}$ term, and constant term $\kappa$.
Consequently the mean of the lattice sum is $\kappa/\alpha$.

For $\ell\ne0$, integrating its Fourier coefficient over $0\leq\theta\leq1$
and unfolding the lattice gives
\[
 \frac1\alpha\left\{
 \Gamma(s_\ell)\zeta(1+s_\ell)-\frac1{s_\ell^2}\right\},
 \qquad s_\ell=\frac{2\pi\ii\ell}{\alpha}.
\]
The coefficient of $\alpha\theta(1-\theta)/2$ is
$-\alpha/(4\pi^2\ell^2)$, which cancels the contribution of
$-1/(\alpha s_\ell^2)$. Its mean is $\alpha/12$.
This proves \eqref{ctp:eq:phase-fourier}.
The gamma factor has exponential decay on the vertical axis; the zeta
factor and its fixed-order derivatives have at most polynomial growth
there, as follows for example from Euler--Maclaurin continuation.
Hence every differentiated Fourier series converges absolutely.
Nonconstancy was proved in \cref{ctp:cor:oscillation}.
\end{proof}

\writenote{The same Mellin kernel and the same constant occur elsewhere
in ProveIt, for a different product. The Fabius-tree compilation
\emph{Geometric Uniform Convolutions and New Frontiers}
\cite{fabiusH}, in its section ``Mellin analysis and the hidden
log-periodic phase'', treats the Laplace product
$H_q(z)=\prod_{k\ge0}(1-\e^{-zq^k})/(zq^k)$, with kernel
$g(x)=\log\bigl((1-\e^{-x})/x\bigr)$, Mellin transform
$-\Gamma(1+s)\zeta(1+s)/s=-\Gamma(s)\zeta(1+s)$, and constant
$c_1=\pi^2/12-\gamma^2/2-\gamma_1$, which is the $\kappa$ of
\eqref{ctp:eq:kappa}. Neither text proves the other's theorem: the
products, the lattices and the normalizations differ, and the
coincidence is the de Bruijn--Mahler mechanism common to both.}

\subsection{The first two bulk orders, with a uniform remainder}
\begin{theorem}[Moving-lattice expansion]\label{ctp:thm:H}
Uniformly in the fractional part of $m$, as $m\to\infty$,
\begin{align}
 \HH(m)={}&\frac\alpha2m^2-
 \left(\frac32+\frac\alpha2\right)m+
 \frac{15}{8}\log m-\frac98-K+\Psi_\alpha(m)
 +\frac{P_1(m)}m+O(m^{-2}),\label{ctp:eq:H-expansion}\\
 P_1(\theta)={}&\frac{27}{16}-\frac32\partial_\alpha\Psi_\alpha(\theta).
 \label{ctp:eq:P1}
\end{align}
There is a complete extension with $\sum_{j=1}^J P_j(m)m^{-j}$ and
remainder $O_J(m^{-J-1})$. The $P_j$ are smooth and one-periodic. The
expansion may be differentiated any fixed number of times if sufficiently
many terms are retained.
\end{theorem}
\begin{proof}
In \eqref{ctp:eq:exact-decomposition}, put $d=k-m$. For a bounded $d$,
\[
 r(m+d)-r(m)=\alpha d-\frac{\beta d}{m}+O((1+|d|)^2m^{-2}).
\]
Taylor expansion of the localized term therefore begins as
\[
 g(\alpha d)-\frac{\beta d}{m}g'(\alpha d)+O\bigl(m^{-2}
 (1+|d|)^L\e^{-c_0|d|}\bigr)
\]
for a fixed $L$; on the positive side an even stronger tail is available.
One-sided values at $d=0$ cause no difficulty because the multiplier $d$
vanishes. The missing terms with $k\leq0$ in the infinite lattice are
exponentially small.

Expansion of the exact bulk term using \eqref{ctp:eq:r-asymptotic} and
\eqref{ctp:eq:S-asymptotic} gives
\begin{align*}
 Mr(m)-S(M)={}&\frac\alpha2m^2-(\beta+\alpha/2)m+
 (\beta/2-b_1)\log m+b_1-K\\
 &+\frac\alpha2\theta(1-\theta)
 +\frac{27/16-(\beta/2)\theta(1-\theta)}m+O(m^{-2}).
\end{align*}
Combining this with the localized sum gives \eqref{ctp:eq:H-expansion}, since
\[
 \partial_\alpha\Psi_\alpha(\theta)
 =\tfrac12\theta(1-\theta)
  +\sum_{j\in\Z}(j-\theta)g'(\alpha(j-\theta)).
\]

Here is a uniform justification of the arbitrary-order procedure.
Split the lattice into $|d|\leq m^{1/3}$ and its complement.
In the first range, expansions of $r(m+d)-r(m)$ and Taylor expansions of
$g$ have remainders bounded by a fixed polynomial in $|d|$ times
$m^{-J-1}\e^{-c_0|d|}$. The signs of $d$ and of the two arguments of $g$
coincide, so Taylor expansion never crosses the kink. Summing preserves
$O_J(m^{-J-1})$. In the complementary range, monotonicity and the positive
lower bound for $r'$ give an exponentially small tail, including the
initial indices far to the left. The same reasoning applies after any
fixed number of derivatives, using more terms at the start.
At integral $m$, the jumps of the one-sided lattice derivatives cancel
exactly with those of the fractional-part bulk polynomials. This can be
seen either from the exact decomposition, by transferring the boundary
index between the two pieces, or coefficient by coefficient in that
identity. The resulting coefficient functions therefore match with all
fixed-order derivatives across the endpoints of a period.
\end{proof}

\subsection{A constructive rule for every coefficient}
\label{ctp:subsec:coefficient-rule}
The preceding proof can be made into a finite symbolic recipe. Write
\begin{equation}\label{ctp:eq:Dj}
 r(m+d)-r(m)=\alpha d+\sum_{j\geq1}D_j(d)m^{-j},
\end{equation}
where
\begin{equation}\label{ctp:eq:Dj-formula}
 D_j(d)=\frac{\beta(-1)^jd^j}{j}
 +\sum_{h=1}^{j-1}b_h(-1)^{j-h}
       \binom{j-1}{h-1}d^{j-h}.
\end{equation}
For example, $D_1(d)=-\beta d$ and
$D_2(d)=\beta d^2/2-b_1d$. The localized contribution to the coefficient
of $m^{-j}$ is
\begin{equation}\label{ctp:eq:Tj}
 T_j(\theta)=\sum_{d\in\Z-\theta}\sum_{\ell=1}^{j}
 \frac{g^{(\ell)}(\alpha d)}{\ell!}
 \sum_{\substack{r_1+\cdots+r_\ell=j\\r_\nu\geq1}}
       \prod_{\nu=1}^{\ell}D_{r_\nu}(d).
\end{equation}
Each coefficient is thus an absolutely convergent lattice sum of explicit
functions. At $d=0$, all $D_j(0)$ vanish, so arbitrary consistent one-sided
values of the derivatives give the same term, namely zero.

Expand $Mr(m)-S(M)$ with $M=m-\theta$, then add \eqref{ctp:eq:Tj}. This
produces $P_j(\theta)$. Bernoulli polynomials supply the bulk coefficients;
finite compositions of the $D_j$ supply the localized coefficients.
This is an all-orders construction, not merely a claim that unspecified
coefficients exist. Fourier transformation can be used afterward to
replace the localized sums by rapidly convergent Fourier series.

\section{Deriving the explicit coefficient expansion}
\label{ctp:sec:explicit-proof}

\subsection{The saddle coordinate and its correction}
Let $t=\e^{-r(m)}$, and define
\begin{equation}\label{ctp:eq:Q}
 Q(m)=tK_1(t)=\frac{\HH'(m)}{r'(m)}.
\end{equation}
Differentiation of \eqref{ctp:eq:H-expansion}, followed by series division,
gives
\begin{align}
 Q(m)&=m+q_0(m)+\frac{q_1(m)}m+O(m^{-2}),\label{ctp:eq:Qexp}\\
 q_0&=-\frac12+\frac{\Psi_\alpha'}\alpha,
 &q_1&=\frac{P_1'}\alpha+\frac{\beta\Psi_\alpha'}{\alpha^2}.
 \label{ctp:eq:q01}
\end{align}
The phase functions on the right are evaluated at $m$.
The scaled variance satisfies the exact identity
\begin{equation}\label{ctp:eq:B2Q}
 B_2=t^2K_2=Q+\frac{Q'}{r'},
\end{equation}
so that
\begin{equation}\label{ctp:eq:B2exp}
 B_2=m+q_0+\frac{1+q_0'}\alpha+O(m^{-1}).
\end{equation}
More generally, the recurrence
\begin{equation}\label{ctp:eq:B-recursion}
 B_{j+1}=jB_j+\frac{B_j'}{r'}
\end{equation}
generates every scaled cumulant from $Q$. This is useful both for the
proof and for the all-orders construction.

At the saddle, $n=\e^{r(m)}Q(m)$. Consequently
\begin{equation}\label{ctp:eq:logn-m}
 \log n=\alpha m-\frac12\log m+c+
 \frac{b_1+q_0(m)}m+
 \frac{b_2+q_1(m)-q_0(m)^2/2}{m^2}+O(m^{-3}).
\end{equation}
The core coordinate \eqref{ctp:eq:mu} removes precisely the first three
terms. Write
\begin{equation}\label{ctp:eq:m-reversion}
 m=\mu+\frac{a(\mu)}\mu+\frac{b(\mu)}{\mu^2}+O(\mu^{-3}).
\end{equation}
Taylor substitution into \eqref{ctp:eq:logn-m} gives the explicit functions
\begin{align}
 a&=-\frac{b_1+q_0}{\alpha},\label{ctp:eq:a}\\
 b&=\frac{(1/2-q_0')a-(b_2+q_1-q_0^2/2)}\alpha.
 \label{ctp:eq:b}
\end{align}
All functions here are evaluated at $\mu$.
The residual derivative in \eqref{ctp:eq:logn-m} is $\alpha+O(1/m)$,
which is bounded away from zero. Thus the formal coefficient comparison
also gives the asserted analytic error by the mean value theorem.
Repeating the argument proves an all-orders reversion with smooth periodic
coefficients.

\subsection{The normalization and the first correction}
Taking logarithms of \eqref{ctp:eq:saddle-expansion}, using
\eqref{ctp:eq:E-universal}, and writing $K_2=\e^{2r(m)}B_2$ gives
\begin{equation}\label{ctp:eq:logp-saddle}
 \log p(n)=\HH(m)+Q(m)-r(m)-\frac12\log(2\pi B_2(m))
 -\frac1{12m}+O(m^{-2}).
\end{equation}
Let $d=b_1+q_0$. Substituting the preceding expansions into this equation
first produces the expression in the \emph{true} saddle coordinate:
\begin{align}
 \log p(n)={}&\frac\alpha2m^2-Am+\eta\log m+C+
 \Psi_\alpha(m)+d(m)+\frac{U(m)}m+O(m^{-2}),\label{ctp:eq:logp-m}\\
 U={}&P_1+q_1-b_1-\frac{q_0}{2}
 -\frac{1+q_0'}{2\alpha}-\frac1{12}.
 \label{ctp:eq:U}
\end{align}
The constant simplifies because
\[
 -K-c-\tfrac12\log(2\pi)
 =\tfrac14\log2+\tfrac34\log\pi+\log G(3/2)=C.
\]

Now substitute \eqref{ctp:eq:m-reversion}. At order one,
$\alpha a+d=0$, so the apparently additional phase derivative in $d$
\emph{cancels}. At order $1/\mu$, the coefficient is
\begin{equation}\label{ctp:eq:D-unsimplified}
 \alpha b+(-A+\Psi_\alpha'+q_0')a+U.
\end{equation}
Using \eqref{ctp:eq:a}--\eqref{ctp:eq:b}, this simplifies first to
\[
 P_1-\frac7{12}-\frac1{2\alpha}
 +(1/2-b_1)q_0-\frac12q_0^2-\frac{q_0'}{2\alpha},
\]
and then, with $b_1=-9/8$ and \eqref{ctp:eq:P1}, exactly to
\eqref{ctp:eq:D}. In particular the $P_1'$ terms disappear. This explicit
cancellation is one reason to distinguish $m$ and $\mu$ throughout.
Equations~\eqref{ctp:eq:main-log} and \eqref{ctp:eq:main-multiplicative} follow.

\subsection{Completion of the all-orders proof}
For a fixed requested order $J$, use the following finite operations.
First compute sufficiently many $b_j$ from \eqref{ctp:eq:bj}, bulk coefficients
from \eqref{ctp:eq:S-exact}, and $P_j$ from \eqref{ctp:eq:Tj}. Next form $Q$ and
$B_2,B_3,\ldots$ by \eqref{ctp:eq:Q} and \eqref{ctp:eq:B-recursion}. Substitute
these in the finite Gaussian polynomials \eqref{ctp:eq:Ej} and take the
logarithm of the saddle prefactor and correction. Finally revert
\eqref{ctp:eq:logn-m} about $\mu$ and substitute.

The derivative-controlled estimates in \cref{ctp:thm:H} justify each finite
calculus operation. The saddle theorem supplies the coefficient-extraction
remainder. The reversion has derivative bounded away from zero, so it
preserves the desired error once one extra order has been retained before
substitution into the quadratic term. All coefficient operations are
addition, multiplication, differentiation, and division by series with
nonzero constant term in an algebra of smooth periodic functions. Thus
no new unbounded coefficient, such as $\log\mu$, is generated beyond the
single displayed $\eta\log\mu$. This proves
\eqref{ctp:eq:all-orders-explicit} and completes \cref{ctp:thm:explicit}.

A safe implementation rule is to retain two orders more than the final
requested order during the $m$-to-$\mu$ substitution. This is a bookkeeping
rule, not a new analytic assumption.

\section{Inverse growth and the integer staircase}
\label{ctp:sec:inverse-proof}

\subsection{Reversion of the quadratic logarithmic core}
Start with $Y=\log y$ and $s=\sqrt{2Y/\alpha}$. The leading quadratic
term suggests $\mu=s+A/\alpha+o(1)$. Set
\[
 \mu=s+\frac A\alpha+\frac{v(s)}s+
 O\!\left(\frac{(\log s)^2}{s^2}\right).
\]
Substitution into \eqref{ctp:eq:main-log}, at constant order after the
quadratic and linear terms cancel, yields
\begin{equation}\label{ctp:eq:v-inverse}
 v(s)=\frac{A^2}{2\alpha^2}
 -\frac{\eta\log s+C+\Psi_\alpha(s+A/\alpha)}\alpha.
\end{equation}
Since $\log n=\alpha\mu-\tfrac12\log\mu-\tfrac12\log\pi$,
\eqref{ctp:eq:inverse} follows by Taylor expansion. The term
$D_\alpha(\mu)/\mu$ in $\log p$ first affects the next inverse order,
not the displayed $1/R$ correction.

Higher inverse terms are obtained by the same procedure. The leading
inverse derivative is of size $s$, so a residual term is removed by
division by $s$ and a bounded perturbation. Every differentiation of a
periodic coefficient leaves it in the same coefficient algebra, while
products of the already occurring $\log s$ terms produce polynomials in
$\log s$. The phase is always reducible by Taylor expansion to
$s+A/\alpha$. This proves the asserted form of the all-orders inverse.

\subsection{Why a continuous calculation describes the discrete threshold}
Adding a unit part proves that $p(n)$ is nondecreasing. Moreover,
$p(n)-p(n-1)$ counts partitions of $n$ without a unit part.
Every integer $n\geq4$ is a nonnegative combination of $2$ and $5$:
even $n$ use only $2$, and odd $n\geq5$ use one $5$ and then $2$.
Therefore $p(n)>p(n-1)$ for all $n\geq4$, and the threshold
\eqref{ctp:eq:threshold} is unambiguous for large $y$.

Let $f_J(x)$ denote the smooth right side of
\eqref{ctp:eq:all-orders-explicit}, with $n$ replaced by a large positive real
$x$. Its derivative with respect to $X=\log x$ satisfies
\begin{equation}\label{ctp:eq:inverse-slope}
 \frac{\dd}{\dd X}f_J(\e^X)=\mu(\e^X)+O(1)>0.
\end{equation}
The asymptotic theorem says, at integers,
$\log p(n)=f_J(n)+O(\mu(n)^{-J-1})$ with a uniform constant.
If $f_J(x)=Y$, a change of $O(\mu^{-J-2})$ in $\log x$ moves $f_J(x)$
by the entire coefficient error. Bracketing the corresponding real
arguments by neighboring integers therefore gives
\begin{equation}\label{ctp:eq:inverse-transport}
 \log N(y)-\log x
 =O(\mu(x)^{-J-2})+O(x^{-1}),
\end{equation}
provided $x$ is the inverse of the chosen smooth truncation. The bracket
uses only monotonicity, not an assumed canonical interpolation of $p$.

The $O(x^{-1})$ term comes from rounding an argument to an integer; it is
smaller than every fixed inverse power of $R$. The first term is the
analytic truncation uncertainty and is much larger at a fixed order.
Accordingly, the theorem gives an asymptotic formula for the threshold,
not an exact integer reconstruction algorithm for every $y$.
Equations~\eqref{ctp:eq:v-inverse} and \eqref{ctp:eq:inverse-transport} prove
\cref{ctp:thm:inverse}.

\subsection{A practical caution}
The expansion parameter is approximately $1/\log n$, not $1/n$.
Even at $n=10^6$, the effective number of active parts is only about
$11.36$. The leading inverse equivalent can therefore have a large
finite-argument error. The first correction helps substantially, but a
numerical application at moderate size should use the implicit saddle
formula and invert it numerically, preferably with a certified remainder
if an integer decision is required. Section~\ref{ctp:sec:computation} reports
actual errors rather than presenting an asymptotic equivalent as a
uniformly accurate approximation.

\section{The largest Catalan part in a random partition}
\label{ctp:sec:largest}
Choose uniformly among the $p(n)$ partitions of $n$. Let $J_n$ be the
largest index $k$ for which a part of size $c_k$ occurs. For $n>0$ this
is well-defined. It is natural to center $J_n$ at the saddle cutoff, not
at $\log_4 n$ alone.

\begin{theorem}[A phase-dependent largest-index law]\label{ctp:thm:largest}
Let $m=m_n$ be the true saddle coordinate, $M=\lfloor m\rfloor$, and
$\theta_n=m-M$. For every fixed integer $h$,
\begin{equation}\label{ctp:eq:largest-law}
 \PP(J_n\leq M+h)=
 \prod_{j=h+1}^{\infty}\left(1-\exp\{-4^{j-\theta_n}\}\right)
 +O_h(m_n^{-1}).
\end{equation}
In particular, along any subsequence with $\theta_n\to\theta\in[0,1)$,
$J_n-M$ converges in distribution to the integer-valued law with cumulative
probabilities
\begin{equation}\label{ctp:eq:largest-cdf}
 G_h(\theta)=\prod_{j=h+1}^{\infty}(1-\e^{-4^{j-\theta}}).
\end{equation}
The estimates are for fixed $h$; a growing-$h$ or moment theorem is not
being asserted.
\end{theorem}
\begin{proof}
Under the product probability measure at parameter $t$, let the numbers
of parts of sizes $c_k$ be independent geometric random variables:
\[
 \PP(Z_k=z)=(1-\e^{-tc_k})\e^{-tc_k z},\qquad z=0,1,2,\ldots.
\]
Conditioning on $\sum_k c_kZ_k=n$ gives the uniform distribution on
partitions of $n$, since every such partition has the same product
weight. This provides an exact probabilistic interpretation, not an
approximation.

Put $K=M+h$ and split the product at $K$. Let $H_\mathrm{pre}$ and
$H_\mathrm{tail}$ be its logarithms. At the full saddle, the removed tail
has scaled cumulants of every fixed order $O_h(1)$. To see this, write
$k=M+j$ and use
\[
 tc_{M+j}=\exp\{r(M+j)-r(m)\},
\]
which is bounded away from zero for each fixed starting $j=h+1$ and
then grows geometrically. The corresponding sums of the functions
$f_r$ from \cref{ctp:lem:cumulants} are bounded. Thus the removed mean is
$O_h(1/t)$, whereas the full standard deviation is $\asymp\sqrt m/t$.

Apply the same coefficient saddle argument to the prefix product. Its
saddle differs from $t$ by $O_h(t/m)$: the mean deficit is $O_h(1/t)$
and the derivative of the mean is of size $m/t^2$. Taylor's theorem
then shows that its saddle action differs from the prefix action
evaluated at $t$ by $O_h(1/m)$. The variance ratio is $1+O_h(1/m)$,
and the two first saddle correction factors are also $1+O_h(1/m)$.
The arc estimates remain uniform because the prefix still contains
$m-O_h(1)$ active parts and all the fixed small factors used in
Section~\ref{ctp:sec:saddle}. Consequently,
\[
 \frac{[z^n]\prod_{k\leq K}(1-z^{c_k})^{-1}}{p(n)}
 =\e^{-H_\mathrm{tail}(t)}\bigl(1+O_h(m^{-1})\bigr)
 =\prod_{k>K}(1-\e^{-tc_k})\bigl(1+O_h(m^{-1})\bigr).
\]
This ratio is exactly the desired conditional probability.

Finally, for each fixed $j$,
\[
 tc_{M+j}=4^{j-\theta_n}
 \left(1+O\left(\frac{1+|j|}{m}\right)\right).
\]
For $0\leq j-\theta_n\leq m/2$, the corresponding logarithmic
error is bounded by $C(1+|j|)/m$. The derivatives of the logarithmic
tail factors have a summable, superexponentially decaying majorant
throughout this range. The range $j-\theta_n>m/2$ contributes a
superexponentially small tail, and the finitely many negative indices
are controlled separately. Thus summation replaces the tail product
by \eqref{ctp:eq:largest-cdf} with error $O_h(m^{-1})$. This proves \eqref{ctp:eq:largest-law}.
For each phase, $G_h\to1$ as $h\to+\infty$ and $G_h\to0$ as
$h\to-\infty$, so it is the cumulative distribution of a proper
integer-valued law. Pointwise convergence at the integer cutoffs gives
the distributional assertion.
\end{proof}

The local scale $4^{j-\theta}$ in this law is the same scale that produces
the Fourier amplitude in \eqref{ctp:eq:phase-fourier}. Both results describe
the bounded-width transition between parts that are typically available
in many copies and parts that are rarely used. The phase is not noise:
it records where the continuous cutoff falls between two consecutive
Catalan indices.

\section{A universality statement for polynomially perturbed geometric parts}
\label{ctp:sec:universality}
The preceding method separates universal local behavior from sequence-specific
normalization. The following proposition makes that separation precise.

\begin{proposition}[General exponential--polynomial part growth]
\label[proposition]{ctp:prop:general}
Suppose an increasing integer part sequence $a_k$ satisfies the sufficient
saddle hypotheses at the end of Section~\ref{ctp:sec:saddle}. Suppose its logarithm
has a smooth eventual interpolation with a termwise differentiable expansion
\[
 r(k)=\alpha k-\beta\log k+c+b_1/k+b_2/k^2+\cdots,\qquad\alpha>0.
\]
Define $K_{\aset}$ by
\begin{align*}
 \sum_{k=1}^M\log a_k={}&\frac\alpha2M^2-\beta M\log M
 +(\alpha/2+c+\beta)M\notag\\
 &+(b_1-\beta/2)\log M+K_{\aset}+O(M^{-1}).
\end{align*}
Let $p_{\aset}(n)$ be the corresponding partition number, and let $\mu$ be
the large solution of $n=\e^c\e^{\alpha\mu}\mu^{1-\beta}$. Then
\begin{align}
 \log p_{\aset}(n)={}&\frac\alpha2\mu^2+
 (1-\beta-3\alpha/2)\mu+
 (3\beta/2-b_1-1/2)\log\mu\notag\\
 &-K_{\aset}-c-\frac12\log(2\pi)+\Psi_\alpha(\mu)+O(\mu^{-1}).
 \label{ctp:eq:general}
\end{align}
There is an all-orders extension with smooth periodic inverse-power
coefficients. The function $\Psi_\alpha$ is exactly
\eqref{ctp:eq:phase-fourier}; it depends on $\alpha$, not on $\beta$ or the
finite initial segment.
\end{proposition}
\begin{proof}
The constant $K_{\aset}$ exists by Euler--Maclaurin summation of the
assumed expansion, or by subtracting enough explicitly summable terms.
The bulk and local argument gives
\[
 \HH(m)=\frac\alpha2m^2-(\beta+\alpha/2)m+
 (\beta/2-b_1)\log m+b_1-K_{\aset}+\Psi_\alpha(m)+O(m^{-1}).
\]
The local limit $r(m+d)-r(m)\to\alpha d$ is unchanged. Also
$Q=m+q_0+O(m^{-1})$, where $q_0=-1/2+\Psi_\alpha'/\alpha$,
independent of $\beta$. Therefore
\[
 \log n=\alpha m+(1-\beta)\log m+c+(b_1+q_0)/m+O(m^{-2}).
\]
Reversion about $\mu$ again gives
$m=\mu-(b_1+q_0)/(\alpha\mu)+O(\mu^{-2})$.
In the logarithm of the saddle estimate, the constant $b_1+q_0$ cancels
against the quadratic-core displacement just as before. The remaining
linear, logarithmic, and constant terms are exactly \eqref{ctp:eq:general}.
The proof of every higher order is the same finite construction with the
assumed coefficients $b_j$ in place of their Catalan values.
\end{proof}

This proposition explains why merely computing the exponential ratio of
part sizes is insufficient. The coefficient $\beta$ changes the linear
term and the size coordinate; $b_1$ changes the power of $\mu$; and the
whole initial sequence contributes to $K_{\aset}$. The local periodic
shape, on the other hand, is determined by the limiting logarithmic gap.

\section{Computation, independent checks, and reproducibility}
\label{ctp:sec:computation}

\subsection{What the supplied programs do}
The package contains an exact integer dynamic program for $p(0),\ldots,p(N)$,
a separate logarithmic-derivative recurrence check, arbitrary-precision
saddle evaluation, independent Fourier and real-lattice evaluations of the
phase, the explicit first-correction formula, and the inverse formula.
The numerical program implements $E_1$ and $E_2$ directly; a rational
multi-index generator lists the Gaussian polynomials through $E_4$ in the
recorded data. The all-orders theorem does not depend on a claim that an
unlimited symbolic implementation is already supplied.

Exact dynamic programming starts with $p(0)=1$ and $p(n)=0$ for $n>0$.
For each allowed part $c\leq N$, in increasing order, replace
$p(n)$ by $p(n)+p(n-c)$ for $n=c,c+1,\ldots,N$. The ascending update
allows repeated copies of $c$. An independent recurrence is
\begin{equation}\label{ctp:eq:exact-recursion}
 np(n)=\sum_{j=1}^n b(j)p(n-j),\qquad
 b(j)=\sum_{c_k\mid j}c_k,
\end{equation}
obtained from $zF'(z)/F(z)$. Agreement of these two algorithms is checked
through $n=300$. The first 20 terms agree with the inspected OEIS entry.
The stored small coefficient file is independently generated; it is not a
redistributed download of another author's table.

The complete run used Python 3.13.5 and mpmath 1.3.0 at 60 decimal digits.
Exact coefficients were generated through $N=10^6$. The Barnes product
identity was checked at $M=0,\ldots,20$. Fourier and real-lattice formulas
for both $\Psi_\alpha$ and its partial $\alpha$-derivative were compared at
10 phases; their maximum recorded discrepancy was less than
$4.8\cdot10^{-59}$. Rational Gaussian algebra gives the first three
constants in \eqref{ctp:eq:E-universal} exactly.

These statements describe what was actually executed. Decimal output and
finite tail cutoffs are not rigorous interval enclosures, and the tests do
not prove the asymptotic estimates. A separate symbolic audit verifies the
normalization and first-correction cancellation using exact algebra.

\subsection{Coefficient accuracy}
Let $P_0$ denote the saddle prefactor in \eqref{ctp:eq:saddle-expansion},
$P_1=P_0(1+E_1)$, and $P_2=P_0(1+E_1+E_2)$ in this subsection only.
The table reports signed relative error, $P_j/p(n)-1$.

\begin{table}[htbp]
\centering
\caption{Exact-cumulant saddle approximations. Errors are not assumed to be monotone.}
\label{ctp:tab:saddle}
\begin{tabular}{rrrrr}
\toprule
$n$ & $m_n$ & $P_0$ error & $P_1$ error & $P_2$ error\\
\midrule
$10^2$ & 4.5337 & $2.126\cdot10^{-2}$ & $-4.169\cdot10^{-3}$ & $-3.016\cdot10^{-3}$\\
$10^3$ & 6.2400 & $1.832\cdot10^{-2}$ & $7.139\cdot10^{-4}$ & $4.696\cdot10^{-4}$\\
$10^4$ & 7.9492 & $1.380\cdot10^{-2}$ & $2.681\cdot10^{-4}$ & $3.448\cdot10^{-5}$\\
$10^5$ & 9.6549 & $9.577\cdot10^{-3}$ & $-3.205\cdot10^{-4}$ & $-1.160\cdot10^{-4}$\\
$10^6$ & 11.3564 & $8.557\cdot10^{-3}$ & $6.471\cdot10^{-5}$ & $4.661\cdot10^{-5}$\\
\bottomrule
\end{tabular}
\end{table}

For reference, some exact values from the computation are
\begin{align*}
 p(100)&=2031,\\
 p(1000)&=18299196,\\
 p(10000)&=8125868781802,\\
 p(100000)&=180507423003606259924,\\
 p(1000000)&=198409222879647727339514603123.
\end{align*}
The explicit $\mu$ formula is more transparent but converges more slowly
at these arguments. Exponentiating the right side of
\eqref{ctp:eq:main-log} without, and with, $D_\alpha/\mu$ gives the following
signed relative errors.

\begin{table}[htbp]
\centering
\caption{Explicit periodic formulas at moderate arguments.}
\label{ctp:tab:explicit}
\begin{tabular}{rrr}
\toprule
$n$ & Through the periodic constant term & Including $D_\alpha/\mu$\\
\midrule
$10^2$ & $-0.122792$ & $-0.069993$\\
$10^3$ & $-0.072947$ & $-0.036697$\\
$10^4$ & $-0.048618$ & $-0.018121$\\
$10^5$ & $-0.040289$ & $-0.012929$\\
$10^6$ & $-0.032146$ & $-0.010559$\\
\bottomrule
\end{tabular}
\end{table}

As a check independent of coefficient extraction, substitute
$m=10.2,30.2,100.2,300.2$ directly into the exact product logarithm and
subtract \eqref{ctp:eq:H-expansion} through $P_1/m$. The residual multiplied
by $m^2$ is respectively about
$-0.1474,-0.2747,-0.3078,-0.3166$. This is consistent with the next
periodic coefficient at the fixed phase $0.2$. It is not used to infer
the remainder theorem.

\subsection{Inverse and largest-index checks}
At $y=p(n)$ and $n\geq4$, the true threshold is exactly $N(y)=n$.
This permits an unambiguous inverse test.

\begin{table}[htbp]
\centering
\caption{Signed relative error of the inverse approximations at $y=p(n)$.}
\label{ctp:tab:inverse}
\begin{tabular}{rrr}
\toprule
$n$ & Leading equivalent \eqref{ctp:eq:inverse-leading} & Corrected formula \eqref{ctp:eq:inverse}\\
\midrule
$10^3$ & $2.042368$ & $0.133978$\\
$10^4$ & $1.548609$ & $0.071730$\\
$10^5$ & $1.256439$ & $0.044236$\\
$10^6$ & $1.061185$ & $0.029094$\\
\bottomrule
\end{tabular}
\end{table}
The large leading errors are displayed deliberately: an asymptotic
equivalent can be mathematically correct while still being a poor
numerical approximation at a million. The corrected formula is much
better in this test, but is not a certified integer inversion rule.

At $n=10^6$ we have $M=11$ and $\theta_n\approx0.35639211$.
Prefix dynamic programs give exact numerators for the event in
\cref{ctp:thm:largest}. The comparison is
\begin{table}[htbp]
\centering
\caption{Largest-index probabilities at $n=10^6$.}
\label{ctp:tab:largest}
\begin{tabular}{rrrr}
\toprule
$h$ & Cutoff $M+h$ & Exact probability & $G_h(\theta_n)$\\
\midrule
$-2$ & 9 & 0.05550439 & 0.05898171\\
$-1$ & 10 & 0.39681499 & 0.41691590\\
$0$ & 11 & 0.89851234 & 0.91283590\\
$1$ & 12 & 0.99999711 & 0.99994242\\
\bottomrule
\end{tabular}
\end{table}
The exact probabilities in the table mean ratios of exactly computed
integers, displayed as decimals. Their limits retain the phase because
only a bounded number of indices occupy the transition region.

\subsection{Numerical safeguards and build instructions}
The cumulant implementation works with the scaled quantities $B_j=t^jK_j$
to avoid unnecessary overflow. It evaluates $1-\e^{-y}$ using
\texttt{-expm1(-y)} near zero. The geometric sums in \eqref{ctp:eq:H-K} are
computed as rational functions using Eulerian polynomials, rather than
by summing a slowly convergent inner series.

For the Fourier coefficients, the supplied code explicitly requests the
Euler--Maclaurin method for $\zeta(1+2\pi\ii\ell/\log4)$ and its derivative.
At even $\ell$, the denominator $1-2^{1-s}$ of an eta-based formula for
$\zeta(s)$ vanishes at the evaluation point, so a naive numerical quotient
is ill-conditioned. In the tested mpmath environment, forcing the
Euler--Maclaurin route removed a discrepancy exposed by the independent
real-lattice check. This is a safeguard for the tested implementation,
not a claim about every current zeta library. The relevant evaluation
methods are documented by mpmath~\cite{mpmath}.

From the package directory, run
\begin{verbatim}
python code/verify.py --max-n 1000000 --dps 60
python code/symbolic_audit.py
pdflatex -interaction=nonstopmode -halt-on-error article.tex
pdflatex -interaction=nonstopmode -halt-on-error article.tex
\end{verbatim}
The CSV files are convenient for inspecting the tables, and
\texttt{data/verification.json} records the environment, exact checks,
constants, and additional Gaussian coefficients. The code performs no
network access. The bibliography and source-audit notes record the external
provenance separately.

\writenote{In the collection these commands must not be run inside the
report directory. \texttt{code/verify.py} writes by default to the
\texttt{data/} directory beside \texttt{code/}, that is to the shipped
data, unless \texttt{--output} names another directory;
\texttt{code/symbolic\_audit.py} always writes
\texttt{../data/symbolic\_audit.json} relative to its own location and
has no option; and \texttt{code/build.sh} changes into \texttt{code/},
where there is no \texttt{article.tex}, so it fails as shipped. Rerun
on a copy of the report directory (the collection README gives the
commands). At intake both scripts were rerun on such a copy: the three
CSV files were byte-identical to the delivery,
\texttt{a033552\_0\_300.txt} and \texttt{symbolic\_audit.json} equal
apart from Windows line ends, and \texttt{verification.json} equal
apart from line ends and the recorded Python and platform strings.}

\section{Further research questions and concrete next steps}
\label{ctp:sec:future}
The following directions are deliberately separated from the proved
statements. Each would improve either the mathematical scope or the
reliability of computation.

\subsection{Certified errors and exact inverse decisions}
Replace the unspecified constants in the contour and moving-lattice
remainders by explicit, computable bounds. The proof already divides the
problem into positive sums, elementary arc inequalities, and Taylor
remainders. An interval implementation could enclose the saddle, truncate
both lattice tails with proven bounds, and return a rigorous interval for
$\log p(n)$. The decisive inverse question is then whether that interval
separates two neighboring integer thresholds. This is substantially
stronger than a fixed-order asymptotic equivalent.

\subsection{A symbolic compiler for periodic asymptotic coefficients}
Implement \eqref{ctp:eq:Tj}, \eqref{ctp:eq:B-recursion}, \eqref{ctp:eq:Ej}, and the
$m$-to-$\mu$ reversion in one coefficient algebra. A useful output would
include $L_2,L_3,L_4$ in a normalized basis of phase derivatives, together
with exact rational audits. The central design issue is to avoid treating
$\Psi_\alpha(m)$ as a constant under differentiation and to track the
partial $\alpha$-derivative separately from the phase derivative.

\subsection{Beyond-all-orders contributions}
The proof discards the compact outer arcs at a relative scale
$O(t\sqrt m)$, much smaller than every inverse power of $m$. Their actual
contributions may have an arithmetic organization at roots of unity.
Identify the dominant secondary arcs, their exponential actions, and how
they interact with the analytic tail of the Barnes expansion. A complete
transseries including such sectors would go beyond the all-orders
Poincar\'e expansion established here.

\subsection{Joint statistics near the largest part}
For finitely many indices $M+j$, the unconditioned geometric variables
suggest a phase-dependent local product limit. Prove a joint conditional
limit with an explicit first correction, rather than only the tail-absence
probability in \cref{ctp:thm:largest}. Extend the estimates uniformly in $h$
far enough to justify moment convergence and asymptotics for the expected
largest index. Neither extension follows automatically from the fixed-$h$
statement.

\subsection{The number of parts and mixed conditioning}
Introduce a marking variable for the total number of parts. Small allowed
parts and the bounded-width cutoff have very different fluctuation scales.
Determine the correct normalization and a joint limit with $J_n$.
A two-variable saddle may need a different treatment from a direct
Gaussian approximation because the unit-part multiplicity can remain
macroscopic on the scale $1/t$.

\subsection{Fuss--Catalan and other gamma-quotient part sequences}
Apply \cref{ctp:prop:general} to integer part sequences with a gamma-quotient
asymptotic, after checking the initial-part gcd condition and removing
any duplicated values. Calculate the normalization constant by Barnes
or multiple-gamma products when possible. The proposition predicts a
universal periodic shape for each limiting logarithmic gap, but the
power and the normalization require a new calculation in every family.

\subsection{Exact counting complexity}
The supplied dynamic program uses $O(n\log n)$ integer additions to
compute all coefficients through $n$, because there are $O(\log n)$
allowed parts up to $n$. That is not a polynomial-in-$\log n$ algorithm
for a single exact value. Determine whether Catalan-part partitions admit
a substantially faster exact counting method of the kind motivated by
Pak's discussion~\cite{pak}. The asymptotic formulas may provide useful
bounds or numerical guidance, but they do not settle this question.

\subsection{Formalization and a periodic coefficient algebra}
The finite coefficient identities, the rational Gaussian multi-index
formula, the gamma-quotient recurrence for Catalan numbers, and the
monotonicity argument for $N(y)$ are natural first formalization targets.
The analytic frontier is a reusable theorem for positive Euler products
with a controlled central arc and fixed-factor minor-arc suppression.
For ProveIt's transseries development, an algebra of smooth or analytic
periodic coefficient germs would encode the phase correctly. One should
not label the present contour proof formalized merely because some
underlying algebraic operations have Lean counterparts.

\section{Conclusion}
A033552 admits a concrete, computable all-orders asymptotic description
whose most visible unusual feature is a persistent periodic amplitude.
The dominant size coordinate is not simply $\log_4 n$: the Catalan
polynomial factor shifts the saddle and changes both the logarithmic
power and the inverse growth. Keeping that shift, the exact Barnes
constant, the local lattice correction, and the Gaussian coefficient
corrections in one calculation yields \eqref{ctp:eq:main-log} and its inverse.

The main analytic issues have explicit proofs: noncentral coefficient
arcs are negligible to every fixed inverse-logarithmic order; local
lattice Taylor expansions have uniform remainders; saddle reversion has
a derivative bounded away from zero; and integer threshold errors are
separated from analytic inverse errors. The largest-index theorem gives
a probabilistic consequence of the same local phase. What remains is not
a missing leading argument, but sharper certification, higher explicit
coefficients, secondary exponential sectors, wider statistics, and a
systematic formalization.

\appendix
\crefalias{section}{appendix}% write (batch 73O1): references print "appendix"
\section{A compact audit of the algebra}
\label{ctp:app:audit}
This appendix records the dependency chain for the constants most likely
to be affected by an indexing or sign error.

\begin{center}
\begin{tabular}{ll}
\toprule
Quantity & Formula or origin\\
\midrule
Allowed unit part & $c_1=1$ only; no second $c_0$ factor\\
Catalan logarithmic drift & $\beta=3/2$, $b_1=-9/8$, $b_2=1/2$\\
Summatory logarithmic coefficient & $b_1-\beta/2=-15/8$\\
Barnes constant & $K=-\tfrac34\log(2\pi)-\log G(3/2)$\\
Bulk $1/m$ coefficient & $27/16$ before the phase correction\\
First local phase correction & $-(3/2)\partial_\alpha\Psi_\alpha$\\
Saddle mean correction & $q_0=-1/2+\Psi_\alpha'/\alpha$\\
Logarithmic power of $\mu$ & $3\beta/2-b_1-1/2=23/8$\\
Normalized constant & $C=-K-c-\tfrac12\log(2\pi)$\\
First Gaussian log correction & $-1/(12m)$\\
\bottomrule
\end{tabular}
\end{center}

The cancellation producing $D_\alpha$ is the identity
\begin{align*}
 &\alpha b+(-A+\Psi_\alpha'+q_0')a+
 P_1+q_1-b_1-\frac{q_0}{2}-\frac{1+q_0'}{2\alpha}-\frac1{12}\\
 &\hspace{8mm}=\frac16-\frac1{2\alpha}
 -\frac32\partial_\alpha\Psi_\alpha
 +\frac{17}{8\alpha}\Psi_\alpha'
 -\frac{(\Psi_\alpha')^2+\Psi_\alpha''}{2\alpha^2},
\end{align*}
with $a,b,q_0,q_1$ as in \eqref{ctp:eq:q01} and
\eqref{ctp:eq:a}--\eqref{ctp:eq:b}. The supplied symbolic audit treats the phase
and its derivatives as independent symbols and verifies this rational
identity exactly. This check isolates coefficient algebra from the
separate analytic justification of the asymptotic operations.

\section{Provenance and scope of the search}
\label{ctp:app:provenance}
The external check was performed on October 1, 2026. It included the OEIS
entry and linked survey, searches for \texttt{A033552 asymptotic} and
\texttt{partitions into Catalan numbers asymptotic}, the classical
periodicity reference, and the ProveIt transseries documentation and
repository code search. The phrase ``Catalan partitions'' alone is
ambiguous in the literature and was not treated as identifying this
sequence without checking the definition.

No claim of an exhaustive MathSciNet, zbMATH, or journal-archive search is
made. The absence of a posted OEIS asymptotic is an observable entry-level
gap, not evidence that no general theorem could already imply a leading
term. Accordingly the main theorems are stated on their mathematical
content, and the article does not claim a resolved named conjecture or
verified historical priority. No repository file or OEIS entry was edited
as part of this work.

\section{Provenance of the collection edition}
\label{ctp:app:edition}

\writenote{This appendix was added in the write.}

\paragraph{Source.} The report is built from one manuscript, number~39
of batch~73O1 of ProveIt's incoming-reports intake: the archive
\texttt{Catalan\_Partitions\_Research.zip} (17~files, a 21-page PDF),
which arrived in commit \texttt{f8c3a392a} and was placed in the
collection by commit \texttt{9df4ba51a}. Its pinned repository commit is
\texttt{e6190a94552f119e9a2140e5fea8450a045a0837} (1~October~2026); the
delivered \texttt{SOURCE\_NOTES.md} also names
\texttt{1f1981f682b2878bde51a6ad40c22777f362fc05} and wrongly calls it a
tree (see the note in \cref{ctp:sec:intro}). The title-page author line
is kept as delivered; the delivered PDF metadata read ``Research report
prepared with ChatGPT'', and the write replaced that field by the
collection's ``AI-assisted research report prepared for Vladimir
Reshetnikov''. The report is AI-assisted and unrefereed.

\paragraph{What the write changed.} Every one of the 93 delivered labels
received the prefix \texttt{ctp:}, and every reference was updated with
it; this appendix adds one label. The write added the
\texttt{[Write note]} paragraphs (after the principal-formula box, in
\cref{ctp:sec:intro} twice, after the proof of \cref{ctp:lem:Fourier},
in \cref{ctp:sec:computation} and here), the bibliography entries
\cite{debruijn} and \cite{fabiusH}, the \verb|\writenote| macro, the PDF
author field above, cleveref type hints (\verb|\label[lemma]|,
\verb|\label[corollary]|, \verb|\label[proposition]|) on the five labels
of lemmas, corollaries and propositions, and
\verb|\crefalias{section}{appendix}| after \verb|\appendix|. These
environments share the theorem counter, and the delivered PDF printed
every reference to them as ``theorem''; only those printed names
changed, no label name or number. No delivered sentence, statement,
proof or table was changed or removed.

\paragraph{Delivered names and shipped paths.} \texttt{article.tex},
\texttt{README.md} (rewritten as the collection's report guide) and
\texttt{SOURCE\_NOTES.md} are at the report root; \texttt{code/} keeps
its three Python files, and \texttt{build.sh} moved there from the
package root; \texttt{data/} keeps its seven files, and
\texttt{requirements.txt} moved there from the package root. All but the
article and the README are byte-identical to the delivery. Not shipped:
the delivered \texttt{article.pdf} (replaced by a build of this
edition) and the checksum ledger \texttt{SHA256SUMS.txt} (verified 16/16
at intake, then retired). \texttt{data/build\_and\_quality\_checks.json}
describes the delivered 21-page PDF.

\paragraph{Relation to repository material.} No repository report or
formal development treats A033552. The Mellin computation of
\cref{ctp:lem:Fourier} reappears, for a different product, in the
Fabius-tree compilation cited in the note after it \cite{fabiusH}. The
batch's other partition reports, on square-root-restricted partitions
(A097356) and on distinct-partition norms (A022629), study different
products in different saddle regimes and share only the standard
saddle, Edgeworth and Lambert-$W$ toolkit; no proposition is proved in
two of them.

\begin{thebibliography}{9}
\bibitem{oeis}
The OEIS Foundation Inc., \emph{The On-Line Encyclopedia of Integer
Sequences}, A033552, ``Number of partitions into Catalan numbers.''
Entry inspected October 1, 2026. \url{https://oeis.org/A033552}.

\bibitem{pak}
I. Pak, \emph{Complexity problems in enumerative combinatorics},
arXiv:1803.06636v2, 2018, especially \S2.5.
\url{https://arxiv.org/abs/1803.06636}.

\bibitem{erdosrichmond}
P. Erd\H{o}s and B. Richmond,
``Concerning periodicity in the asymptotic behaviour of partition functions,''
\emph{Journal of the Australian Mathematical Society, Series A}
\textbf{21} (1976), 447--456.
DOI: \href{https://doi.org/10.1017/S1446788700019285}{10.1017/S1446788700019285}.
Author archive: \url{https://www.renyi.hu/~p_erdos/1976-07.pdf}.

\bibitem{fgd}
P. Flajolet, X. Gourdon, and P. Dumas,
``Mellin transforms and asymptotics: Harmonic sums,''
\emph{Theoretical Computer Science} \textbf{144} (1995), 3--58.
\url{https://specfun.inria.fr/dumas/Publications/FlGoDu95.pdf}.

\bibitem{dlmfG}
NIST Digital Library of Mathematical Functions, \S5.17,
``Barnes' $G$-function (double gamma function).''
\url{https://dlmf.nist.gov/5.17}. Accessed October 1, 2026.

\bibitem{dlmfzeros}
NIST Digital Library of Mathematical Functions, \S25.10(i),
``Zeros: Distribution.''
\url{https://dlmf.nist.gov/25.10.i}. Accessed October 1, 2026.

\bibitem{proveit}
V. Reshetnikov, \emph{ProveIt}, repository snapshot
\texttt{e6190a94552f119e9a2140e5fea8450a045a0837};
\path{Analysis/Transseries/README.md} and
\path{Analysis/Transseries/docs/series-and-transseries/README.md}.
\url{https://github.com/VladimirReshetnikov/ProveIt/tree/e6190a94552f119e9a2140e5fea8450a045a0837/Analysis/Transseries}.

\bibitem{mpmath}
The mpmath developers, \emph{Zeta functions, L-series and polylogarithms},
mpmath documentation, including Euler--Maclaurin evaluation options.
\url{https://mpmath.org/doc/1.1.0/functions/zeta.html}.
The accompanying numerical run used mpmath 1.3.0; see the recorded
version in \texttt{data/verification.json}.

\bibitem{debruijn}
[Added in the write.] N.~G. de Bruijn, ``On Mahler's partition
problem,'' \emph{Indagationes Mathematicae} \textbf{10} (1948),
210--220.

\bibitem{fabiusH}
[Added in the write.] ProveIt, Fabius-tree compilation
\emph{Geometric Uniform Convolutions and New Frontiers},
\path{Analysis/FabiusFunction/docs/semi-formalized-research-frontiers/drafts/frontier-compilations/Geometric_Uniform_Convolutions_and_New_Frontiers/fabius-frontier-report-H.tex},
section ``Mellin analysis and the hidden log-periodic phase'' (Lemma
``Mellin kernel'').
\end{thebibliography}
\end{document}
